% =============================================================================
% FINAL REPORT — Time-Series-Enhanced Predictive Process Monitoring
% MSc Project 1 – SPI Lab, RWTH Aachen
% =============================================================================
\documentclass[runningheads]{llncs}

\usepackage[T1]{fontenc}
\usepackage{graphicx}
\usepackage{float}
\usepackage{booktabs}
\usepackage{tabularx}
\usepackage{color}
\usepackage[colorlinks=true, linkcolor=black, citecolor=black, urlcolor=blue]{hyperref}
\renewcommand\UrlFont{\color{blue}\rmfamily}
\urlstyle{rm}
\usepackage{soul}
\usepackage{subcaption}
\usepackage{amsmath}
\usepackage{multirow}

\begin{document}

\title{MSc Project 1: Time-Series-Enhanced Predictive Process Monitoring}
\subtitle{Final Report}
\titlerunning{Time-Series-Enhanced PPM: Final Report}

\author{Micha\u{l} Durkalec \and
    Konstantinos Sakkas \and
    Roman Le\u{z}ovi\u{c}
}
\authorrunning{M. Durkalec et al.}

\institute{Supervised Process Intelligence (SPI) Lab,\\
    RWTH Aachen University, Aachen, Germany\\
    \email{office@pads.rwth-aachen.de}}

\maketitle

% =============================================================================
\begin{abstract}

Predictive process monitoring uses historical event log data to predict
outcomes and remaining times of running process instances, with standard
approaches deriving features from the event log alone while ignoring
temporal resource-load variation. This project tests whether aligned
time-series features (active-case counts with rolling-window statistics)
improve prediction quality. We built a modular four-stage pipeline for the
BPI Challenge 2017 loan application log, comparing log-only Random Forest
models against log+time-series models for binary outcome classification
(approved vs. not-approved) and remaining-time regression. The hourly-resampled
time-series features with 8-hour rolling windows added no meaningful improvement
for classification (F1-weighted +0.1\%) and degraded regression
(RMSE +1.45). We
analyse why this extraction strategy fails to improve either task and
discuss implications for feature engineering in predictive process monitoring
using our open-source and fully reproducible pipeline.

\keywords{Predictive Process Monitoring \and Process Mining \and Time-Series Features
    \and Random Forest \and Model Evaluation \and Explainability}
\end{abstract}

% =============================================================================
\section{Introduction}
\label{sec:introduction}

Predictive process monitoring (PPM) predicts future states of running process
instances from event log data, with common tasks including binary
outcome prediction (will this case end normally or deviate?) and remaining-time
regression (how many hours until completion?). Most PPM approaches derive
features from the event log only, such as activity counts, elapsed time, and case
attributes, capturing what happened in a case but not the
concurrent workload on the process. The BPI Challenge 2017
dataset~\cite{bpic2017} records a loan application process at a Dutch
financial institution with validation, offer, and approval stages, where
resource contention during peak periods can affect both
outcomes and processing times. Log-based features alone do not capture
this workload variation over time.

This project answers: \emph{do aligned time-series workload features
improve predictive performance for outcome and remaining-time tasks compared
to event-log features alone?} We built a modular pipeline that extracts
log-based features, computes time-series features, trains multiple different models,
and evaluates them against baselines. The contributions are: an open-source, reproducible PPM
pipeline with time-series support, an empirical comparison of log-only vs.
log+time-series models on both tasks, and an analysis of when time-series
features help and when they do not.

The paper is organised as follows. \autoref{sec:data} describes the dataset and
preprocessing. \autoref{sec:features} covers feature engineering.
\autoref{sec:models} specifies the models and evaluation protocol.
\autoref{sec:results} presents results. \autoref{sec:testing} describes testing
evidence. \autoref{sec:evaluation} discusses findings and limitations.
\autoref{sec:conclusion} concludes.

% =============================================================================
\section{Data and Preprocessing}
\label{sec:data}

\subsection{Event Log: BPI Challenge 2017}
\label{sec:bpic2017}
The BPI Challenge 2017 event log~\cite{bpic2017} records a loan application
process at a Dutch financial institution, with cases describing the lifecycle
of credit applications including application processing, offer creation,
document validation, and customer interactions. Events belong to three
categories: application state changes, offer state changes, and workflow
activities performed by employees.

The dataset is widely used in predictive process monitoring as it contains rich control-flow, temporal, and financial information along with substantial variation in case duration and outcome. \autoref{tab:data-stats} summarises the dataset.

\begin{table}[H]
    \centering
    \caption{BPI Challenge 2017 dataset statistics.}
    \label{tab:data-stats}
    \begin{tabularx}{0.6\textwidth}{lX}
        \toprule
        Statistic & Value \\
        \midrule
        Cases & 31,509 \\
        Events & 1,202,267 \\
        Activities & 26 \\
        Start date & 2016-01-01 \\
        End date & 2017-02-01 \\
        Mean case length (events) & 38.16 \\
        Median case length (events) & 35 \\
        \bottomrule
    \end{tabularx}
\end{table}

\subsection{Prediction Tasks}
% TODO [RL] Why was my reference removed? Where is the references.bib?
This study evaluates model performance on two predictive process monitoring tasks: loan outcome classification and remaining-time regression. Cases terminate as approved, declined, or cancelled, corresponding to the
occurrence of \texttt{A\_Pending}, \texttt{A\_Denied}, or
\texttt{A\_Cancelled} respectively. For outcome prediction, we formulate a binary classification task with approved cases as the positive class while declined and cancelled cases form a single negative class. For remaining-time prediction, the target variable is the time in hours between the last event of a prefix and case completion.

\subsection{Preprocessing Pipeline}
\label{sec:preprocessing}

Raw XES event logs are first ingested, cleaned, and sorted in ascending order by timestamp before a sliding-window approach generates training instances. For each case, a window of length $k$ at time step $t$ consists of the most recent $k$ events ending at position $t$, with $k$ constrained by a minimum and maximum window length. Each resulting window is associated with the corresponding case label. Table~\ref{tab:preprocessing-summary} summarises the final dataset sizes for each prediction task along with the applied filtering steps and sliding-window configuration. 

The preprocessing pipeline differs between the regression and classification tasks. For regression, no additional filtering is applied as sliding windows are generated with a minimum length of 3 and a maximum length of 100 events per case, producing truncated event histories that capture recent temporal context. For classification, additional constraints ensure label consistency and prevent information leakage: cases without any defined outcome events are removed (98 out of 31,509 cases), and windows containing an outcome event are excluded since their inclusion would leak information about the final case outcome into the input representation. The same minimum and maximum window lengths (3 to 100 events) are used as in the regression setup.

\begin{table}[H]
    \centering
    \caption{Preprocessing summary per task.}
    \label{tab:preprocessing-summary}
    \begin{tabularx}{\textwidth}{lX}
        \toprule
        Component & Value \\
        \midrule

        Initial dataset &
        31,509 cases / 1,202,267 events \\

        Window type &
        Sliding history window \\

        Window length &
        min = 3, max = 100 \\

        Regression filtering &
        None (all cases retained) \\

        Classification filtering &
        No-outcome cases removed (98 excluded), \\ & outcome prefixes removed \\

        Regression final size &
        31,509 cases \\

        Classification final size &
        31,411 cases \\

        \bottomrule
    \end{tabularx}
\end{table}

\subsection{Temporal Train-Test Split}
\label{sec:split}

We use a temporal split by case start time to prevent leakage, ensuring all
prefixes from the same case remain in the same split. Cases starting before the 0.8
quantile go to training while later cases go to testing, simulating an
operational setting where the model predicts on cases that start after the
training period. \autoref{tab:split-stats} shows the resulting split
statistics.

\begin{table}[H]
    \centering
    \caption{Temporal split statistics per prediction task.}
    \label{tab:split-stats}
    \begin{tabularx}{\textwidth}{lXX}
        \toprule
        & Train & Test \\
        \midrule

        \multicolumn{3}{l}{\textbf{Classification task}} \\
        Cases & \texttt{22,859} & \texttt{6,283} \\
        Prefixes & \texttt{763,417} & \texttt{210,096} \\
        Split timestamp & \multicolumn{2}{X}{\texttt{2016-10-18 11:39:41}} \\
        
        \midrule
        \multicolumn{3}{l}{\textbf{Regression task}} \\
        Cases & \texttt{22,895} & \texttt{6,302} \\
        Prefixes & \texttt{827,614} & \texttt{228,455} \\
        Split timestamp & \multicolumn{2}{X}{\texttt{2016-10-18 22:39:59}} \\

        \bottomrule
    \end{tabularx}
\end{table}

% =============================================================================
\section{Feature Engineering}
\label{sec:features}

\subsection{Log-Based Features (Baseline)}
\label{sec:log-features}

Log-based features are derived purely from the event log, with \autoref{tab:log-features}
listing the feature groups. Activity counts are one-hot encoded per event type
while temporal features use the timestamp of the last event in the prefix, with all
features computed in vectorised form for efficiency.

\begin{table}[H]
    \centering
    \caption{Log-based feature groups.}
    \label{tab:log-features}
    \begin{tabularx}{\textwidth}{lXr}
        \toprule
        Group & Features & Count \\
        \midrule
        Activity counts & Count per event type per prefix & 26 \\
        Temporal & Elapsed time, time since last event & 2 \\
        Financial & Monthly cost, number of terms & 2 \\
        Waiting states & Time in waiting activities & 4 \\
        \midrule
        Total & & 34 \\
        \bottomrule
    \end{tabularx}
\end{table}

\subsection{Time-Series Features (Enhanced)}
\label{sec:ts-features}

Time-series features capture concurrent process workload at hourly
resolution, where for every hourly base signal we compute three features: the
raw value, an 8-hour trailing mean (\texttt{rolling\_mean\_8h}), and a
trend (\texttt{trend\_8h} = current value minus value 8 hours earlier).
\autoref{tab:ts-features} lists the time-series feature classes.

\begin{table}[H]
    \centering
    \caption{Time-series feature classes.}
    \label{tab:ts-features}
    \begin{tabularx}{\textwidth}{lXrr}
        \toprule
        Class & Base signal & Base cols & Total \\
        \midrule
        ActiveCaseCount & Cases active at prefix timestamp & 1 & 3 \\
        FinancialVolume & Hourly sum/mean of withdrawal, offer, cost & 6 & 18 \\
        DecisionRate & Ratio of accepted decisions per hour & 1 & 3 \\
        \midrule
        Total & & 8 & 24 \\
        \bottomrule
    \end{tabularx}
\end{table}

\subsection{Feature Extraction Pipeline}
\label{sec:extraction}

Features are extracted via a modular pipeline where each feature set implements the
\texttt{PrefixFeature} protocol (\texttt{name()}, \texttt{fit()},
\texttt{\_\_call\_\_()}, \texttt{feature\_names}), and the feature extractor
generates a single feature matrix by concatenating all enabled feature sets.
Time-series features are precomputed once during \texttt{fit()} and aligned
to each prefix at extraction time via vectorised timestamp lookup.

% =============================================================================
\section{Model Specification and Evaluation Protocol}
\label{sec:models}

\subsection{Prediction Targets}
\label{sec:tasks}
\autoref{tab:class-distribution} shows the class distributions for the train and test splits at both trace and prefix level, with the distributions highly consistent between sets and indicating a stable split without noticeable sampling bias. At trace level, a moderate class imbalance towards approved cases can be observed, becoming more pronounced at prefix level where approved cases form a larger proportion of the data.

\begin{table}[H]
    \centering
    \caption{Class distribution for traces and prefixes (train/test split).}
    \label{tab:class-distribution}
    \begin{tabularx}{0.65\textwidth}{l l l r c}
        \toprule
        Level & Split & & Count & Percentage (\%) \\
        \midrule
        \multirow{4}{*}{Traces}
         & \multirow{2}{*}{Train} & approved & 12,694 & 55.5\% \\
         & & not approved & 10,165 & 44.5\% \\
         \cmidrule(lr){2-5}
         & \multirow{2}{*}{Test} & approved & 3,449 & 54.9\% \\
         & & not approved & 2,834 & 45.1\% \\
        \midrule
        \multirow{4}{*}{Prefixes}
         & \multirow{2}{*}{Train} & approved & 516,099 & 67.6\% \\
         & & not approved & 247,318 & 32.4\%  \\
         \cmidrule(lr){2-5}
         & \multirow{2}{*}{Test} & approved & 141,171 & 67.2\% \\
         & & not approved & 68,925 & 32.8\% \\
         \bottomrule
    \end{tabularx}
\end{table}


\autoref{tab:remaining-time-distribution} shows remaining-time statistics at prefix level, with the target distribution highly right-skewed as indicated by the substantial difference between mean and median values in both train and test sets. The mean remaining time is approximately 300 hours while the median is considerably lower (around 205 hours), indicating that many prefixes occur relatively late in the process where little time remains until case completion, whereas a smaller number of early prefixes contribute disproportionately large remaining-time values. This leads to a long upper tail in the distribution with maximum values exceeding 3000 hours in the training set, while the interquartile range further reflects this variability with 50\% of prefix instances lying approximately between 40 and 480 hours remaining. 

The train and test distributions are highly consistent across all statistics, suggesting that the temporal structure of prefixes is preserved across the split with no significant distribution shift.

\begin{table}[H]
    \centering
    \caption{Distribution of remaining time (in hours) for train and test prefix sets.}
    \label{tab:remaining-time-distribution}
    \begin{tabular}{l r r}
        \toprule
        Statistic & Train & Test \\
        \midrule
        Prefixes & 827,614 & 228,455 \\
        Mean     & 302.8 & 300.0 \\
        Median   & 205.1 & 208.8 \\
        Std      & 315.5 & 298.1 \\
        Min      & 0   & 0 \\
        Max      & 3269.0 & 2348.1 \\
        Q1       & 42.3  & 45.3 \\
        Q3       & 479.4 & 488.8 \\
        \bottomrule
    \end{tabular}
\end{table}



\subsection{Models and Hyperparameter Optimisation}
\label{sec:model-list}

We evaluate four model families for both prediction tasks (\autoref{tab:models}),
covering linear baselines and non-linear ensemble methods. The selected models
are widely used in predictive process monitoring due to their ability to handle
high-dimensional, sparse, and heterogeneous event-based feature representations.

% TOOD [RL] Are thee correct? Final runconfigs don't do parameter search? 
\begin{table}[H]
    \centering
    \caption{Models and hyperparameter search spaces.}
    \label{tab:models}
    \begin{tabularx}{\textwidth}{llX}
        \toprule
        Model & Task & Key hyperparameters \\
        \midrule
        Random Forest & Both &
            \texttt{n\_estimators}: 100--300\newline
            \texttt{max\_depth}: 5--20\newline
            \texttt{min\_samples\_split}: 2--20 \\
        Logistic Regression & Classification & \texttt{C}: 0.01--100 \\
        Ridge Regression & Regression & \texttt{alpha}: 0.01--100 \\
        HistGradientBoosting & Both &
            \texttt{learning\_rate}: 0.01--0.3\newline
            \texttt{max\_iter}: 100--300 \\
        \bottomrule
    \end{tabularx}
\end{table}

Hyperparameter optimisation is performed using \texttt{RandomizedSearchCV}
with 3-fold cross-validation on the training set. All models
are trained exclusively on training prefixes, and evaluation is performed on a
held-out test set.

\autoref{tab:hpo-config} lists the search configuration.

\begin{table}[H]
    \centering
    \caption{Hyperparameter optimisation configuration.}
    \label{tab:hpo-config}
    \begin{tabularx}{0.7\textwidth}{lX}
        \toprule
        Parameter & Value \\
        \midrule
        Search method & RandomizedSearchCV \\
        CV folds & 3 \\
        Search iterations & 10 \\
        Search sample size & All training prefixes \\
        Scoring & F1-weighted (classification) / RMSE (regression) \\ % TODO [RL] We don't do this in our runconfigs, do we?
        \bottomrule
    \end{tabularx}
\end{table}

\subsection{Evaluation Protocol and Comparison Design}
\label{sec:metrics}

We select three metrics for classification. \textbf{F1-weighted} balances
precision and recall while accounting for class support, making it robust to
the 67:33 class imbalance as the primary ranking metric used by the
pipeline. \textbf{ROC AUC} provides a threshold-independent measure of
discrimination as the most commonly reported metric in PPM literature.
\textbf{MCC} (Matthews Correlation Coefficient) aggregates all four confusion
matrix entries into a single score (range $-1$ to $+1$, where 0 = random) and
acts as a balanced tiebreaker when F1 and ROC AUC disagree.

For regression we select \textbf{RMSE} in hours (same unit as the target,
penalises large errors, primary ranking metric) and \textbf{$R^2$}
(proportion of variance explained, normalised against the mean baseline).

Every
model configuration is trained twice: once with log-only features and once
with log-based and time-series features combined. Both configurations share
the same train/test split, hyperparameter search space, and evaluation
pipeline. The comparison tool (\texttt{evaluate\_feature\_impact.py}) loads
both evaluation reports and produces paired metric tables and faceted plots.

% =============================================================================
\section{Results}
\label{sec:results}

The best model differs by task: HistGradientBoosting (HGB) achieves the
highest F1-weighted for classification (0.6601), while Random Forest (RF)
achieves the lowest RMSE for regression (228.97).

\subsection{Outcome Prediction}
\label{sec:results-outcome}

% TODO [RL] Waht is the featureset here? Full feature set? Are we using time series features here?
% TODO [RL] How are these scores computed? Over all prefixes? per prefix then aggregated. Explain how these were computed explicitly.
HistGradientBoosting (HGB) outperforms all alternatives for outcome
classification, achieving the highest scores across all three primary
metrics (\autoref{tab:classification-results}). HGB leads Random Forest
by 2.1\% in F1-weighted (0.6601 vs.\ 0.6465), by 0.006 in ROC AUC, and
by 0.019 in MCC.

\begin{table}[H]
    \centering
    \caption{Classification model comparison (log+TS, 58 features, full data).}
    \label{tab:classification-results}
    \begin{tabularx}{0.85\textwidth}{lccc}
        \toprule
        Model & F1-weighted & ROC AUC & MCC \\
        \midrule
        HistGradientBoosting & \textbf{0.6601} & \textbf{0.6580} & \textbf{0.2316} \\
        Random Forest        & 0.6465 & 0.6522 & 0.2128 \\
        Logistic Regression  & 0.6297 & 0.6270 & 0.1827 \\
        Dummy (most frequent)& 0.5401 & 0.5000 & 0.0000 \\
        \bottomrule
    \end{tabularx}
\end{table}

\autoref{fig:classification-prefix} shows F1-weighted across prefix lengths,
with performance improving consistently as more events become available: HGB
rises from 0.5051 at prefix~3 to 0.7731 at prefix~34. The improvement rate
diminishes after the first 10--15 events, but longer prefixes continue to add signal.

\begin{figure}[H]
    \centering
    \includegraphics[width=\textwidth]{figures/f1_weighted_vs_prefix.png}
    \caption{F1-weighted by prefix length (classification).}
    \label{fig:classification-prefix}
\end{figure}

\paragraph{Time-Series Impact}
Adding hourly-resampled TS features to HGB changes nothing. With 34 log-only
features, HGB scores 0.6584 F1-weighted and 0.6597 ROC AUC. With 58
log+TS features, it scores 0.6594 and 0.6557 (\autoref{tab:classification-ts}).
The differences (+0.001 in F1-weighted and $-$0.004 in ROC AUC) are
negligible.

\begin{table}[H]
    \centering
    \caption{Classification TS impact (HGB, full data).}
    \label{tab:classification-ts}
    \begin{tabularx}{0.85\textwidth}{lccc}
        \toprule
        Configuration & Features & F1-weighted & ROC AUC \\
        \midrule
        HGB log-only (34) & 34 & 0.6584 & 0.6597 \\
        HGB log+TS (58)   & 58 & 0.6594 & 0.6557 \\
        \bottomrule
    \end{tabularx}
\end{table}

\paragraph{Feature Importance}
\autoref{tab:classification-feature-importance} lists the top 10
features for HGB log-only. A single feature dominates:
\texttt{count\_W\_Validate application} (importance 0.0636) is three
times the second-ranked feature and accounts for the majority of
predictive signal.

\begin{table}[H]
    \centering
    \caption{Top 10 features by permutation importance (HGB log-only).}
    \label{tab:classification-feature-importance}
    \begin{tabularx}{0.85\textwidth}{lr}
        \toprule
        Feature & Importance \\
        \midrule
        count\_W\_Validate application        & 0.0636 \\
        count\_W\_Handle leads               & 0.0209 \\
        monthly\_cost                         & 0.0123 \\
        num\_terms                            & 0.0097 \\
        elapsed\_time\_hours                  & 0.0074 \\
        count\_A\_Validating                  & 0.0061 \\
        count\_O\_Returned                    & 0.0060 \\
        count\_W\_Call after offers           & 0.0057 \\
        count\_W\_Call incomplete files       & 0.0043 \\
        time\_since\_last\_event\_hours       & 0.0023 \\
        \bottomrule
    \end{tabularx}
\end{table}

\autoref{fig:classification-feature-importance} visualises the
importance distribution, with activity-count features occupying 6 of the top
10 positions while financial and temporal features appear at middle ranks
with substantially lower importance and waiting-state features appearing
nowhere in the top 10.

\begin{figure}[H]
    \centering
    \includegraphics[width=\textwidth]{figures/hist_gradient_boosting_feature_importance.png}
    \caption{Permutation feature importance (HGB log-only).}
    \label{fig:classification-feature-importance}
\end{figure}

\autoref{fig:shap-classification} shows the SHAP summary dot plot
for the same model, where high occurrences of
\texttt{count\_W\_Validate application} push predictions toward
approved while higher \texttt{elapsed\_time\_hours} and
\texttt{time\_since\_last\_event\_hours} push toward
not-approved, meaning long-running cases end negatively.

\begin{figure}[H]
    \centering
    \includegraphics[width=\textwidth]{figures/shap_hgb_dot.png}
    \caption{SHAP summary dot plot (HGB classification). Each dot represents a
             single prediction. Colour indicates feature value (red = high,
             blue = low). Features are ordered by global importance.}
    \label{fig:shap-classification}
\end{figure}

\subsection{Remaining-Time Prediction}
\label{sec:results-time}

Random Forest (RF) achieves the lowest RMSE for remaining-time regression,
marginally ahead of HGB and well below the dummy baseline
(\autoref{tab:regression-results}). The $R^2$ of 0.382 indicates that
RF explains 38\% of the variance in remaining time.

\begin{table}[H]
    \centering
    \caption{Regression model comparison (log+TS, max\_length=100, full data).}
    \label{tab:regression-results}
    \begin{tabularx}{0.75\textwidth}{lrr}
        \toprule
        Model & RMSE (hours) & $R^2$ \\
        \midrule
        Random Forest        & \textbf{230.28} & 0.382 \\
        HistGradientBoosting & 231.15  & -- \\
        Ridge                & 237.33  & -- \\
        Dummy (mean)         & 295.97  & 0.000 \\
        \bottomrule
    \end{tabularx}
\end{table}

\autoref{fig:regression-prefix} shows RMSE across prefix lengths as it
improves from 284.83 at prefix~3 to 233.14 at prefix~100, with most of the
gain achieved by prefix~20. Beyond that, returns diminish but do not
fully plateau.

\begin{figure}[H]
    \centering
    \includegraphics[width=\textwidth]{figures/rmse_vs_prefix.png}
    \caption{RMSE by prefix length (RF log-only regression).}
    \label{fig:regression-prefix}
\end{figure}

\paragraph{Time-Series Impact}
Time-series features degrade regression performance. The log-only RF
(228.97 RMSE) beats every log+TS variant by 1.31--1.45 RMSE
(\autoref{tab:regression-ts}), a 0.6\% penalty. The degradation is consistent across all
prefix lengths.

\begin{table}[H]
    \centering
    \caption{Regression TS impact (RF, max\_length=100, full data).}
    \label{tab:regression-ts}
    \begin{tabularx}{0.75\textwidth}{lrr}
        \toprule
        Configuration & Features & RMSE (hours) \\
        \midrule
        RF log-only (34) & 34 & \textbf{228.97} \\
        RF log+TS (58)   & 58 & 230.42 \\
        RF log+TS (compare) & 58 & 230.28 \\
        \bottomrule
    \end{tabularx}
\end{table}

\paragraph{Full-Length vs.\ Maximum-Length Prefixes}
Capped (\texttt{max\_length=100}) and unbounded (full trace) prefix
configurations produce nearly identical results at 230.28 vs.\ 230.20 RMSE,
with a difference of 0.04\% indicating that longer histories do not help.

\paragraph{Feature Importance}
\autoref{tab:regression-feature-importance} lists the top 10 features.
Top is \texttt{count\_W\_Validate application} (0.0618), then
\texttt{count\_W\_Call after offers} (0.0323) and
\texttt{count\_A\_Cancelled} (0.0299). Activity counts fill
7 of the top 10. Temporal features are less important.

\begin{table}[H]
    \centering
    \caption{Top 10 features by permutation importance (RF log-only).}
    \label{tab:regression-feature-importance}
    \begin{tabularx}{0.85\textwidth}{lr}
        \toprule
        Feature & Importance \\
        \midrule
        count\_W\_Validate application      & 0.0618 \\
        count\_W\_Call after offers         & 0.0323 \\
        count\_A\_Cancelled                 & 0.0299 \\
        count\_A\_Validating                & 0.0132 \\
        elapsed\_time\_hours                & 0.0131 \\
        count\_O\_Returned                  & 0.0127 \\
        count\_A\_Complete                  & 0.0111 \\
        time\_since\_last\_event\_hours     & 0.0096 \\
        count\_O\_Accepted                  & 0.0068 \\
        num\_terms                          & 0.0048 \\
        \bottomrule
    \end{tabularx}
\end{table}

\subsection{Overfitting Analysis}
\label{sec:overfitting}

\autoref{tab:overfitting} reports train-test gaps for the best models while
\autoref{fig:train-vs-test} visualises the comparison.

\begin{table}[H]
    \centering
    \caption{Train-test performance gap.}
    \label{tab:overfitting}
    \begin{tabularx}{\textwidth}{llrr}
        \toprule
        Task & Model & Train & Test (gap) \\
        \midrule
        Classification & HGB (f1-weighted) & 0.6971 & 0.6609 ($-$0.0362) \\
        Classification & RF (f1-weighted)   & 0.6810 & 0.6510 ($-$0.0300) \\
        Regression     & RF log-only (RMSE) & 248.08 & 233.14 ($-$14.94) \\
        \bottomrule
    \end{tabularx}
\end{table}

\begin{figure}[H]
    \centering
    \includegraphics[width=0.7\textwidth]{figures/train_vs_test.png}
    \caption{Train vs.\ test performance (all models).}
    \label{fig:train-vs-test}
\end{figure}

Overfitting is moderate with the train-test gap for HGB at $-$0.036
F1-weighted and for RF regression at $-$14.94 RMSE (6\% of the
test RMSE), both within expected bounds for tree-based ensemble models
without explicit regularisation.

\subsection{Summary of Findings}
\label{sec:results-summary}

\begin{itemize}
    \item \textbf{HGB is best for classification} (0.6601 F1-weighted,
          0.6580 ROC AUC, 0.2316 MCC), outperforming RF by 2.1\%.
    \item \textbf{RF is best for regression} (228.97 RMSE log-only,
          R$^2$=0.382), marginally ahead of HGB.
    \item \textbf{Time-series features add no value} for
          classification (+0.1\% F1-weighted) and \textbf{harm}
          regression (+1.45~RMSE).
    \item \textbf{Log-only RF (228.97~RMSE)} beats all log+TS
          variants, including the HPO-tuned compare run (230.28).
    \item \textbf{Unbounded prefixes yield no improvement}
          ($-$0.04\%~RMSE).
    \item \textbf{\texttt{count\_W\_Validate}} is the single
          most predictive feature across both tasks
          (importance 0.0636 classification, 0.0618 regression).
    \item Performance improves with longer prefixes, but
          diminishes after 10--15~events.
    \item Overfitting gaps are moderate (HGB $-$0.036,
          RF $-$14.94~RMSE).
\end{itemize}

% =============================================================================
\section{Testing Evidence}
\label{sec:testing}

The project includes 133 unit and integration tests across 12
test modules:
\begin{itemize}
    \item \textbf{Targets (9)}: Remaining time and outcome label
          construction on synthetic event logs.
    \item \textbf{Preprocessing (10)}: Sliding-window and
          outcome-window generation, temporal split integrity.
    \item \textbf{Log-based features (13)}: Activity counts,
          temporal, financial, and waiting-state features
          with shape, naming, and edge cases.
    \item \textbf{Time-series features (2)}: Extraction and
          alignment to prefixes.
    \item \textbf{Dataset (7)}: Download, extraction, caching,
          and HTTP error handling (1 marked integration).
    \item \textbf{Metrics (33)}: MAE, RMSE, F1, AUC-ROC,
          R\textsuperscript{2}, plateau detection, and model
          comparison correctness.
    \item \textbf{Plots (11)}: Task detection, metric selection,
          and figure generation.
    \item \textbf{SHAP (5)}: Classification and regression
          explanations, graceful skip for non-tree models.
    \item \textbf{Trainer (18)}: Model fitting, artifacts,
          PCA pipeline, and hyperparameter search smoke tests.
    \item \textbf{Config (21)}: Schema validation, estimator
          construction, YAML round-trip, and error handling.
\end{itemize}
All tests pass with a single command as follows:
\begin{verbatim}
poetry run pytest
\end{verbatim}

% =============================================================================
\section{Critical Evaluation}
\label{sec:evaluation}

\subsection{Do Time-Series Features Improve Predictive Performance?}
\label{sec:discussion-value}

For classification, the hourly-resampled time-series features with 8-hour
rolling windows added no meaningful improvement (F1-weighted +0.1\%).
For regression, they consistently degraded performance (RMSE +1.45),
with the log-only model outperforming all log+TS variants across all
metric dimensions.

Several factors may explain the limited impact. First, the time-series
features were never independently validated as carrying a workload signal:
we compute them but do not verify that active-case counts, financial
volumes, or decision rates actually track known busy periods. The null
result may therefore reflect a poor operationalisation of workload rather
than a property of time-series information per se.

Second, a possible \textbf{proxy effect}: log-based features (activity
counts, elapsed time, case-level financial attributes) may already contain
indirect workload information. The dominance of
\texttt{count\_W\_Validate application}, a gate activity, across both tasks
is consistent with this hypothesis, although it could equally reflect
case maturity rather than workload.

Third, several design choices in the extraction template may limit its
effectiveness. The \textbf{hourly resolution} is coarse: the 8-hour
rolling trend (\texttt{trend\_8h}) is computed as a first-order
difference, which amplifies stochastic noise rather than measuring
directional change. The alignment step uses backward
\texttt{merge\_asof}, assigning each prefix the most recent past
workload value, even though the predictive question is forward-looking
as it asks what load the case will experience next. Standard time-series features such
as recent lags (t-1h, t-2h) and seasonal lags (t-24h, t-168h) were
also omitted despite strong priors for daily and weekly cycles in
loan application processing.

Fourth, the \textbf{regression degradation} (+1.45 RMSE) may partly
reflect the curse of dimensionality: adding 24 features to a 34-feature
set increases the search space for tree-based models, and irrelevant
features are known to degrade ensemble performance.

\subsection{Limitations}
\label{sec:discussion-limits}

Several limitations affect the interpretation of these results. All three
time-series signal types
(active-case count, financial volume, decision rate) share a single extraction
template with hourly binning, 8-hour rolling mean, and 8-hour trend, while
other extraction strategies (per-resource load, sub-hourly resolution,
seasonal lags) remain unexplored. The 8-hour window covers only 2.6\% of
the mean case duration (~302 hours), with windows at 24, 72, or 168 hours
potentially capturing workload seasonality better. The pipeline evaluates Random Forest and
HistGradientBoosting as tree-based baselines, since deep learning models
(LSTMs, transformers) may integrate continuous time-series information
differently than axis-aligned tree splits allow. The
temporal split is deterministic, and time-series cross-validation could
provide a more rigorous evaluation. No runtime comparison between the
log-only and log+TS pipelines was performed, which would inform
practitioners of the engineering cost.

All classifiers achieved MCC values below 0.25, indicating modest
performance overall. The primary ranking metric F1-weighted is sensitive
to class prevalence (67:33 imbalance). Permutation importance, used for
feature ranking, is known to be biased when features are correlated, as
the time-series rolling features are mechanically correlated with their raw
signals and with log-based temporal features, which may affect the
importance rankings.

Regarding threats to validity: identical fixed parameters were used for
both log-only and log+TS configurations, mitigating concerns that
hyperparameter search could favour one feature set. The pipeline is
deterministic given fixed random seeds and configuration, ensuring
reliability.

% =============================================================================
\section{Conclusion}
\label{sec:conclusion}

\subsection{Summary}
\label{sec:summary}

This project investigated whether aligned time-series workload features
improve predictive process monitoring. We built a modular four-stage
pipeline for the BPI Challenge 2017 loan application log, extracting
34 log-based features and 24 hourly-resampled time-series features
with 8-hour rolling windows, while comparing four model families per
task. HistGradientBoosting achieved the best classification performance
(0.6601 F1-weighted, 0.6580 ROC AUC, 0.2316 MCC), with Random Forest
achieving the lowest regression error (228.97 RMSE log-only).

The time-series features added no meaningful value for classification
(+0.1\% F1-weighted) and degraded regression (+1.45~RMSE), with the
log-only RF beating all log+TS variants. The gate activity
\texttt{count\_W\_Validate} dominated feature importance across both
tasks (0.0636 classification, 0.0618 regression), while unbounded prefix
lengths provided no improvement over a cap of 100~events. Overfitting
was moderate for both best models.

One possible explanation is a proxy effect, where log-based features may
already capture workload variation indirectly through activity counts.
The findings suggest that practitioners can prioritise log-based features,
particularly activity counts of gate events, without investing in
time-series feature engineering. The pipeline and all experiments are
open-source and fully reproducible.

\subsection*{Data and Code Availability}

\textbf{Source code:}
\url{https://github.com/durki007/spi-time-series}

\textbf{Dataset:}
The BPI Challenge 2017 log is downloaded automatically on first run.

\textbf{Reproduction:}
\begin{enumerate}
    \item \texttt{pip install -r requirements.txt}
    \item \texttt{python -m spi\_time\_series.main}\\
          \texttt{experiments/classification/classification\_compare.yaml}\\
          \texttt{-o results/classification/classification\_compare/}
    \item \texttt{python -m spi\_time\_series.main}\\
          \texttt{experiments/regression/regression\_rf\_log.yaml}\\
          \texttt{-o results/regression/regression\_rf\_log/}
\end{enumerate}

% =============================================================================
% \begin{thebibliography}{8}

%     \bibitem{bpic2017}
%     van Dongen, B.F.: BPI Challenge 2017. Eindhoven University of Technology
%     (2017). \url{https://research.tue.nl/en/datasets/bpi-challenge-2017}

%     \bibitem{ppm-survey}
%     Di Francescomarino, C., Ghidini, C., Maggi, F.M., Milani, F.: Predictive
%     process monitoring methods: Which one suits you best? In: BPM, pp.
%     462--479 (2018)

%     \bibitem{active-cases}
%     B\"anziger, T., Basin, D., Klaise, J., Srivatsa, S.: Active case
%     count feature for predictive process monitoring. In: ICPM, pp. 1--8
%     (2021)

%     \bibitem{breiman-rf}
%     Breiman, L.: Random forests. Machine Learning 45(1), 5--32 (2001)

%     \bibitem{shap}
%     Lundberg, S.M., Lee, S.-I.: A unified approach to interpreting model
%     predictions. In: NeurIPS, pp. 4765--4774 (2017)

%     \bibitem{logistic-regression}
%     Hosmer, D.W., Lemeshow, S., Sturdivant, R.X.: Applied Logistic
%     Regression, 3rd edn. Wiley (2013)

% \end{thebibliography}
\bibliography{refs}
\bibliographystyle{plain}
\appendix
\appendix

\section{Appendix: Additional Results}
\label{app:additional-results}

\subsection{Classification: Full Model Comparison}

\begin{table}[H]
    \centering
    \caption{Classification: full metrics for all models (log+TS, 58 features).}
    \label{tab:app-classification}
    \begin{tabularx}{0.85\textwidth}{lccc}
        \toprule
        Model & F1-weighted & ROC AUC & MCC \\
        \midrule
        HGB & \textbf{0.6601} & \textbf{0.6580} & \textbf{0.2316} \\
        RF  & 0.6465 & 0.6522 & 0.2128 \\
        LR  & 0.6297 & 0.6270 & 0.1827 \\
        Dummy & 0.5401 & 0.5000 & 0.0000 \\
        \bottomrule
    \end{tabularx}
\end{table}

\subsection{Regression: Full Model Comparison}

\begin{table}[H]
    \centering
    \caption{Regression: full metrics for all models (log+TS, max\_length=100).}
    \label{tab:app-regression}
    \begin{tabularx}{0.75\textwidth}{lrr}
        \toprule
        Model & RMSE & $R^2$ \\
        \midrule
        RF  & \textbf{230.28} & 0.382 \\
        HGB & 231.15 & -- \\
        Ridge & 237.33 & -- \\
        Dummy & 295.97 & 0.000 \\
        \bottomrule
    \end{tabularx}
\end{table}

\subsection{TS Impact Summary}

\begin{table}[H]
    \centering
    \caption{TS impact summary across all configurations.}
    \label{tab:app-ts-impact}
    \begin{tabularx}{\textwidth}{llrrr}
        \toprule
        Task & Model & Features & F1 / RMSE & ROC AUC / $R^2$ \\
        \midrule
        Classification & HGB log-only & 34 & 0.6584 & 0.6597 \\
        Classification & HGB log+TS & 58 & 0.6594 & 0.6557 \\
        \midrule
        Regression & RF log-only & 34 & \textbf{228.97} & 0.388 \\
        Regression & RF log+TS & 58 & 230.42 & 0.381 \\
        Regression & RF log+TS (full) & 58 & 230.20 & 0.382 \\
        \bottomrule
    \end{tabularx}
\end{table}

\subsection{Overfitting Details}

\begin{table}[H]
    \centering
    \caption{Train-test comparison for all models.}
    \label{tab:app-overfitting}
    \begin{tabularx}{\textwidth}{llrrr}
        \toprule
        Task & Model & Metric & Train & Test \\
        \midrule
        Classification & HGB & F1-weighted & 0.6971 & 0.6609 \\
        Classification & RF & F1-weighted & 0.6810 & 0.6510 \\
        Classification & LR & F1-weighted & 0.6563 & 0.6388 \\
        Classification & Dummy & F1-weighted & 0.5454 & 0.5401 \\
        \midrule
        Regression & RF & RMSE & 248.08 & 233.14 \\
        Regression & HGB & RMSE & 238.74 & 234.88 \\
        Regression & Ridge & RMSE & 256.93 & 240.87 \\
        Regression & Dummy & RMSE & 295.97 & 295.97 \\
        \bottomrule
    \end{tabularx}
\end{table}

\section{Appendix: Configuration Files}
\label{app:configs}
Experiments are YAML files under
\texttt{experiments/\{task\}/}. Key parameters:
\texttt{data.split\_quantile=0.8},
\texttt{prefix.min\_length=3},
\texttt{search.cv\_folds=3}, and
\texttt{features.enabled\_features} selects
log-based or time-series feature sets.

\section{Appendix: Feature Descriptions}
\label{app:features}

\subsection{Log-Based Features (34 features)}

\begin{itemize}
    \item \textbf{ActivityCountFeatures (26)}: One count per activity\\
          \quad(e.g., \texttt{count\_W\_Validate application},
          \texttt{count\_A\_Accepted}) in the prefix.
    \item \textbf{TemporalFeatures (2)}:
          \texttt{elapsed\_time\_hours} (hours since case start) and
          \texttt{time\_since\_last\_event\_hours} (hours since the
          most recent event).
    \item \textbf{FinancialFeatures (2)}:
          \texttt{monthly\_cost} and \texttt{num\_terms}
          from the loan application's offer.
    \item \textbf{WaitingStateFeatures (4)}:\\
          \texttt{hours\_since\_offer\_sent},
          \texttt{hours\_since\_incomplete},\\
          \texttt{is\_waiting\_for\_customer},
          \texttt{is\_waiting\_for\_bank}.
\end{itemize}

\subsection{Time-Series Features (24 features)}

Three columns per signal: raw, 8-hour trailing mean,
and 8-hour trend (current minus 8~hours prior).

\begin{itemize}
    \item \textbf{ActiveCaseCountFeature (1 base, 3 total)}:
          Number of cases active at the prefix's timestamp in the hourly
          active-case time series.
    \item \textbf{FinancialVolumeFeature (6 base, 18 total)}:
          Hourly aggregates: total and mean of withdrawal amount,
          offer amount, and monthly cost.
    \item \textbf{DecisionRateFeature (1 base, 3 total)}:
          Ratio of accepted decisions (events \texttt{O\_Accepted})
          per hour.
\end{itemize}

\end{document}
